\begin{abstract}

Deep research agents---multi-agent systems that search, read, and write long, citation-grounded reports---are bottlenecked by LLM inference. We present the first content-level characterization of deep research inference at the LLM-call layer (800 tasks, 15,296 calls). We found a defining property: \textbf{content recycling}. Content is generated once, yet later calls serve it again and again. They re-read it during prefill and re-generate it during decode. Recycling is pervasive but heterogeneous: some calls copy their inputs almost verbatim; others rewrite and condense (copy rates 8--91\%, bimodal). \textbf{Once generated, many times served}. Existing techniques cannot collect this waste: (1) cross-request KV reuse perturbs downstream decisions in multi-agent pipelines and degrades report quality; (2) uniform speculative decoding wins on copy-heavy calls but loses on rewrite-heavy ones. We exploit recycling on both sides of inference. On the prefill side, we reuse upstream KV caches and preserve downstream decisions via attention-guided selective recomputation. On the decode side, we route speculative decoding by predicted copy rate: the prediction selects the drafter and sets the draft depth. On mainstream deep research benchmarks, our system cuts end-to-end latency by up to 56.3\% and TTFT by \tbd\% versus vLLM. Report quality is statistically unchanged.

\end{abstract}